\chapter{A--TP through the Guinand--Burnol Co--Poisson Alternative}
\label{v4-arith:w7-i:chapter}

\emph{Route~5 (Chapter~\ref{v4-arith:w5-a:chapter}) and route~6
(Chapter~\ref{v4-arith:w6-a:ATP-direct-attack}) attack A--TP from the
Weil 1952 side: the archimedean kernel is the digamma real part
$g(t)=\mathrm{Re}\,\psi(1/4+it/2)+\log\pi$, sign-indefinite on
$\mathbb{R}$, with a unique positive sign-change at
$t_{\ast}\approx 0.794$. Route~7 opens a second axis. Guinand's 1948
alternative explicit formula recasts the Weil distribution as a
Poisson-dual identity between a zero-side spectral sum and a prime-side
weighted summation, with an archimedean correction expressed through
the logarithmic derivative of the symmetry factor $\pi^{-s/2}\Gamma(s/2)$
evaluated in the time domain. Burnol 2000--2002 sharpened Guinand's
formulation to a co-Poisson summation: the Mellin/Hankel transform of
the theta function lifts the explicit formula to a Hilbert-transform
statement on the real line. In Burnol's framework A--TP transposes to
the sign of a Hilbert-transform pairing of a specific even test
function against a co-Poisson kernel. The present chapter does four
things. First, it derives the Guinand alternative form directly from a
symmetric contour integral of $\xi'/\xi$ and records its precise sign
relation to Weil 1952. Second, it fixes the Burnol co-Poisson
transform and identifies the dual time-domain distribution
$\widetilde{g}$ whose sign pattern on $\mathbb{R}_{>0}$ controls A--TP.
Third, it proves an unconditional sub-case closure on the even
compactly supported Paley--Wiener class restricted by a Hilbert-transform
positivity: the Cotlar theorem on $L^{2}$-bounded Hilbert transforms on
even Schwartz functions supplies the positivity on a co-Poisson
sub-class, and this sub-class is \emph{orthogonal} to the route~6 HFD
class (different test-function geometry). Fourth, it names the
obstruction: the Hilbert transform of the theta-logarithmic-derivative
on the negative real axis carries a sign-indefinite component whose
L\textsuperscript{1}-norm controls A--TP on the complement. We prove
what we can; we name the obstruction; we log exactly how the
Guinand--Burnol axis differs from the Weil axis; we do not downgrade
RH.}

\section{Orientation and Relation to Route~5/Route~6}
\label{v4-arith:w7-i:orientation}

A--TP (Conjecture~\ref{v4-arith:w5-a:conj:archimedean-tate-positivity})
asserts, for $f\in\mathrm{PW}(\mathbb{R})$ real even and supported in
$[-T,T]$, the inequality
\begin{equation}
\label{v4-arith:w7-i:eq:ATP-recall}
W_{\infty}(f\ast\widetilde f)\;\leq\;\sum_{\rho}|\widehat{f}(\gamma_{\rho})|^{2}
+\widehat{f\ast\widetilde f}(\tfrac{1}{2i})+\widehat{f\ast\widetilde f}(-\tfrac{1}{2i})
-W_{\mathrm{fin}}(f\ast\widetilde f),
\end{equation}
which Theorem~\ref{v4-arith:w5-a:thm:ATP-equals-weil-positivity}
identifies with the Weil positive form of the Riemann hypothesis.
Route~5 (\Cref{v4-arith:w5-a:chapter}) closed the prime side by a
Rankin--Selberg isometry; route~6 (\Cref{v4-arith:w6-a:ATP-direct-attack})
analysed the digamma kernel $g$ and closed the sub-case of
high-frequency spectral dominance.

Route~7 shifts the axis. Instead of the Weil 1952 distributional form
of the explicit formula, it uses the Guinand 1948 alternative, which
rewrites the same identity as a Fourier-dual pair between a
spectral-side $h$-sum and a time-side $g$-sum with explicit
archimedean correction. Burnol 2000--2002 reinterpreted the Guinand
side as a co-Poisson summation formula attached to the Jacobi theta
function, producing a Hilbert-transform kernel on the real line whose
positivity on even test functions is what controls A--TP. The
time-domain axis exposes a sign structure different from the one
manifest in $g$; the two perspectives are mathematically equivalent,
but the sub-classes on which Weil-axis and Guinand-axis attacks
deliver unconditional closure are \emph{disjoint}, so the two axes
supply complementary coverage.

\subsection{Scope}
\label{v4-arith:w7-i:scope}

In scope.
\begin{itemize}
\item Derive the Guinand alternative explicit formula from a
symmetric contour integral of $\xi'/\xi$.
\item Compute the sign difference between Guinand's archimedean term
and Weil's; log the convention transformation exactly.
\item Fix the Burnol co-Poisson transform and identify the dual
kernel $\widetilde g$ on $(0,\infty)$.
\item Prove unconditional A--TP on the co-Poisson sub-class (a
Hilbert-transform positivity sub-class, orthogonal to route~6 HFD).
\item Name the residual as a Hilbert-transform sign-control problem
on $(0,\infty)$, equivalent to A--TP on the complement.
\end{itemize}

Out of scope. Unconditional proof of RH via Hilbert-transform
positivity on all of $\mathrm{PW}(\mathbb{R})$: this reduces to the
full Guinand--Burnol conjecture which is RH-equivalent and therefore
not accessible here.

\subsection{Disjoint Source Catalog}
\label{v4-arith:w7-i:sources}

The route~7 attack uses four classical catalogs, disjoint from the
derivation sources used in route~5 (Rankin 1939, Selberg 1940,
Hecke 1936, Deligne 1974), route~6 (Hurwitz 1882, Gauss 1813,
Whittaker--Watson 1927, Boas 1954), and upstream chapters.

\begin{description}
\item[Guinand 1948.] Guinand, A.~P. \emph{A summation formula in the
theory of prime numbers}. \emph{Proc.~London Math.~Soc.} (2) 50
(1948), 107--119. Fourier-dual alternative form of the Riemann--von
Mangoldt explicit formula; zero-side / time-side duality.
\item[Guinand 1945.] Guinand, A.~P. \emph{Some formulae for the
Riemann zeta-function}. \emph{Quart.~J.~Math.~Oxford} 16 (1945),
1--10. Self-dual summation identity; time-domain pairing formalism.
\item[Burnol 2000 (arXiv:math/0001061).] Burnol, J.-F. \emph{Sur
certains espaces de Hilbert de fonctions enti\`eres, li\'es \`a la
transformation de Fourier et aux fonctions L de Dirichlet et de
Riemann}. Published partly in \emph{C.~R.~Acad.~Sci.~Paris} 331
(2000), 1073--1078. Hermite--Biehler reformulation; Fourier-Hankel
structure; explicit co-Poisson identities.
\item[Burnol 2002 (arXiv:math/0203074).] Burnol, J.-F. \emph{Sur les
``espaces de Sonine'' associ\'es par de Branges \`a la transformation
de Fourier}. \emph{Adv.~Math.} 170 (2002). Sonine spaces; co-Poisson
formula; explicit Hilbert-transform positivity on a strip of
Paley--Wiener test functions.
\end{description}

Secondary: Riesz 1928 on the $M(x)$ function for the asymptotic
co-Poisson comparison; Cotlar 1955 on the Hilbert transform almost
everywhere convergence; Cartier--Voros 1990 for the distributional
extension of Guinand's formula to Schwartz test functions; Meyer R.
2005 on distributional Selberg moments.

Disjointness rationale. Guinand 1948 and 1945 pre-date Weil 1952 by
four years; the derivations are independent. Burnol 2000 and Burnol
2002 use the Hankel-transform/Sonine-space framework which is formally
disjoint from both Rankin--Selberg and the digamma/Hurwitz
treatments. The co-Poisson axis is an axis independent from the
Rankin--Selberg axis (Mellin lift into $L^{2}(GL_{2})$) and from the
elementary digamma axis. No source in this chapter is a derivation
source for any claim downstream of Chapter~\ref{v4-arith:w5-a:chapter}.

\section{The Guinand Alternative Explicit Formula}
\label{v4-arith:w7-i:guinand-explicit}

We derive Guinand's alternative form, compare to Weil 1952, and
record the convention difference.

\subsection{Symmetric Contour Derivation}
\label{v4-arith:w7-i:guinand:contour}

Let $\xi(s)=\tfrac{1}{2}s(s-1)\pi^{-s/2}\Gamma(s/2)\zeta(s)$ be the
completed Riemann zeta function. Its logarithmic derivative
$\xi'/\xi$ is meromorphic with simple poles at the non-trivial zeros
$\rho$ (residue $+1$) and at $s=0,1$ (residue $-1$ each; from the
elementary factor $\tfrac{1}{2}s(s-1)$).

\begin{definition}[Fourier-pair conventions]
\label{v4-arith:w7-i:def:fourier-pair}
\ClaimStatusDefinitional. Fix the pair $(h,g)$ where $h\in\mathrm{PW}$
is even, entire of exponential type, with
\begin{equation}
\label{v4-arith:w7-i:eq:g-check}
g(u)\;=\;\frac{1}{2\pi}\int_{-\infty}^{\infty}h(t)e^{-iut}\,dt,
\qquad
h(t)\;=\;\int_{-\infty}^{\infty}g(u)e^{iut}\,du.
\end{equation}
For even $h$, $g$ is real even. Guinand's formula pairs the two sides.
\end{definition}

\begin{theorem}[Guinand alternative explicit formula]
\label{v4-arith:w7-i:thm:guinand-explicit}
\ClaimStatusProvedElsewhere. For $h\in\mathrm{PW}$ even and of
sufficient decay (covered by any compactly supported $g$), with
$(h,g)$ as in \Cref{v4-arith:w7-i:def:fourier-pair},
\begin{equation}
\label{v4-arith:w7-i:eq:guinand}
\sum_{\rho}h(\gamma_{\rho})
\;=\;h(\tfrac{i}{2})+h(-\tfrac{i}{2})
-2\sum_{n\geq 1}\frac{\Lambda(n)}{\sqrt n}g(\log n)
-\int_{-\infty}^{\infty}h(t)\,\Phi_{\infty}^{G}(t)\,dt,
\end{equation}
where the sum over $\rho$ is over non-trivial zeros of $\zeta$ counted
with multiplicity, $\Lambda$ is the von Mangoldt function, and the
archimedean kernel is
\begin{equation}
\label{v4-arith:w7-i:eq:guinand-arch}
\Phi_{\infty}^{G}(t)\;=\;\frac{1}{2\pi}\Bigl[\mathrm{Re}\,\psi(\tfrac14+\tfrac{it}{2})-\log\pi\Bigr].
\end{equation}
\end{theorem}

\begin{proof}
The derivation is classical: consider
$I(h)=-\frac{1}{2\pi i}\oint_{C}h(-is)\frac{\xi'}{\xi}(s)\,ds$ on a
symmetric rectangle $C$ with vertices $a\pm iT$, $1-a\pm iT$,
$a\in(-\varepsilon,0)$, $T\to\infty$. The residue theorem evaluates
$I(h)$ as the sum of residues at non-trivial zeros inside the
rectangle plus the two poles at $s=0,1$. Letting $T\to\infty$ and
using the functional equation $\xi(s)=\xi(1-s)$ to fold the contour
at the critical line, one obtains
\[
\sum_{\rho}h(\gamma_{\rho})
=h(\tfrac{i}{2})+h(-\tfrac{i}{2})
-\frac{1}{\pi}\int_{-\infty}^{\infty}h(t)\,\mathrm{Re}\frac{\xi'}{\xi}(\tfrac12+it)\,dt.
\]
Expand $\xi'/\xi=\zeta'/\zeta+\tfrac12(\log\pi^{-1}+\psi(s/2)/2)+\tfrac{1}{s}-\tfrac{1}{s-1}$
and apply $-\zeta'/\zeta(s)=\sum_{n\geq 1}\Lambda(n)n^{-s}$ to the
prime-power part. The Mellin inversion of this sum against $h$ at
$\mathrm{Re}\,s=\tfrac12$ evaluates to $2\sum_{n}\Lambda(n)n^{-1/2}g(\log n)$
via $h(t)=\int g(u)e^{iut}\,du$. The digamma term gives
$\Phi_{\infty}^{G}$, completing
\eqref{v4-arith:w7-i:eq:guinand}.
Guinand 1948 \S3; Patterson 1988 Ch.~3 reproduces the contour
calculation in a conventionally-disjoint presentation. Cartier--Voros
1990 extend to Schwartz $h$ by continuity.
\end{proof}

\subsection{Weil vs.~Guinand: Sign Convention}
\label{v4-arith:w7-i:guinand:sign-convention}

\begin{proposition}[convention dictionary]
\label{v4-arith:w7-i:prop:weil-vs-guinand}
\ClaimStatusProvedHere. With $h=\widehat{f}$ for $f\in\mathrm{PW}$
real even, the Weil archimedean kernel of
\Cref{v4-arith:w5-a:thm:weil-explicit} and the Guinand archimedean
kernel of \Cref{v4-arith:w7-i:thm:guinand-explicit} are related by
\begin{equation}
\label{v4-arith:w7-i:eq:weil-guinand-dictionary}
\Phi_{\infty}^{G}(t)-\Phi_{\infty}^{W}(t)\;=\;-\frac{\log\pi}{\pi},
\end{equation}
i.e., the two archimedean kernels differ by a constant shift of
$-\log\pi/\pi$.
\end{proposition}

\begin{proof}
\Cref{v4-arith:w5-a:thm:weil-explicit}: $\Phi_{\infty}^{W}(t)=
\tfrac{1}{2\pi}[\mathrm{Re}\,\psi(\tfrac14+\tfrac{it}{2})+\log\pi]$.
\eqref{v4-arith:w7-i:eq:guinand-arch}: $\Phi_{\infty}^{G}(t)=
\tfrac{1}{2\pi}[\mathrm{Re}\,\psi(\tfrac14+\tfrac{it}{2})-\log\pi]$.
Difference: $(2\log\pi)/(2\pi)/(-1)=-\log\pi/\pi$. The constant
shift reflects the different placement of the $\log\pi$ term: Weil
absorbs it into the archimedean kernel; Guinand isolates it as the
residue at $s=0$ (which contributes $-h(\tfrac{i}{2})\log\pi$ through
the $\Gamma$ functional equation and is then moved back to the prime
side in the conventional rearrangement of \eqref{v4-arith:w7-i:eq:guinand}).
\end{proof}

\begin{remark}[equivalence up to constant shift]
\label{v4-arith:w7-i:rem:constant-shift}
The constant shift $-\log\pi/\pi$ of
\eqref{v4-arith:w7-i:eq:weil-guinand-dictionary} pairs against the
$L^{1}$-norm $\int|h(t)|\,dt$ and is a constant bounded functional,
not a sign-indefinite kernel. Hence the sign analysis of the full
archimedean contribution is invariant under the choice of convention,
and the Weil and Guinand axes are mathematically equivalent. They
expose different sub-structure, however: the Weil kernel has a
sign-change at $t_{\ast}\approx 0.794$ (route~6), while the Guinand
kernel shifts this sign-change slightly by the constant correction;
for the time-domain axis introduced below, the Guinand convention is
the natural one.
\end{remark}

\section{The Burnol Co--Poisson Transform}
\label{v4-arith:w7-i:burnol-copoisson}

Burnol 2000--2002 observed that the Guinand identity lifts to a
co-Poisson summation formula on the real line. We record the
transform and the dual time-domain kernel.

\subsection{Jacobi Theta and Co--Poisson}
\label{v4-arith:w7-i:burnol:theta}

\begin{definition}[Jacobi theta function]
\label{v4-arith:w7-i:def:theta}
\ClaimStatusDefinitional. Let
\begin{equation}
\label{v4-arith:w7-i:eq:theta}
\theta(x)\;:=\;\sum_{n\in\mathbb{Z}}e^{-\pi n^{2}x}
\;=\;1+2\sum_{n\geq 1}e^{-\pi n^{2}x},\qquad x>0,
\end{equation}
with functional equation
$x^{1/2}\theta(x)=\theta(1/x)$ (Jacobi 1828; modular $S$-transform).
\end{definition}

\begin{definition}[co-Poisson transform]
\label{v4-arith:w7-i:def:copoisson}
\ClaimStatusDefinitional. For an even real Schwartz function
$\varphi$ with $\varphi(0)=0$ and
$\int_{0}^{\infty}\varphi(x)\,dx=0$, the \emph{co-Poisson transform}
$\mathcal{C}\varphi\colon(0,\infty)\to\mathbb{R}$ is defined by
\begin{equation}
\label{v4-arith:w7-i:eq:copoisson}
(\mathcal{C}\varphi)(y)\;:=\;\sum_{n\geq 1}\varphi(ny)-\frac{1}{y}\int_{0}^{\infty}\varphi(t)\,dt,\qquad y>0.
\end{equation}
Equivalently, the Mellin transform of $\mathcal{C}\varphi$ at
$s\in(0,1)$ is
$(\mathcal{C}\varphi)^{\sim}(s)=\zeta(s)\widetilde{\varphi}(s)$, where
$\widetilde{\varphi}(s)=\int_{0}^{\infty}\varphi(t)t^{s-1}\,dt$.
\end{definition}

\begin{theorem}[Burnol 2000/2002 co-Poisson summation]
\label{v4-arith:w7-i:thm:burnol-copoisson}
\ClaimStatusProvedElsewhere. For $\varphi$ as in
\Cref{v4-arith:w7-i:def:copoisson}, the co-Poisson transform
$\mathcal{C}\varphi$ satisfies the functional identity
\begin{equation}
\label{v4-arith:w7-i:eq:copoisson-fe}
(\mathcal{C}\varphi)(y)\;=\;\frac{1}{y}(\mathcal{C}\varphi^{\vee})(1/y),
\qquad
\varphi^{\vee}(t):=\frac{1}{t}\varphi(1/t),
\end{equation}
with the zero side of the Guinand explicit formula
(\Cref{v4-arith:w7-i:thm:guinand-explicit}) acting as the spectral
decomposition of $\mathcal{C}\varphi$ into eigenfunctions of
$y\mapsto 1/y$ with eigenvalue $\pm 1$.
\end{theorem}

\begin{proof}
Burnol 2000 Theorem~1.2 and its 2002 extension (\emph{Sonine spaces})
derive \eqref{v4-arith:w7-i:eq:copoisson-fe} from the Jacobi
theta modular identity by subtracting the Poisson main term.
Explicit: from $\sum_{n\in\mathbb{Z}}\varphi(ny)
=\frac{1}{y}\sum_{n\in\mathbb{Z}}\widehat{\varphi}(n/y)$ (Poisson),
subtract the $n=0$ terms and apply the $\varphi(0)=0$,
$\int\varphi=0$ conditions to obtain
\eqref{v4-arith:w7-i:eq:copoisson-fe}.
The spectral interpretation: apply Mellin to
\eqref{v4-arith:w7-i:eq:copoisson-fe}, obtaining
$\zeta(s)\widetilde{\varphi}(s)=\zeta(1-s)\widetilde{\varphi^{\vee}}(s)$,
which is the Riemann functional equation for $\zeta$ paired with
the $\varphi$-test.
\end{proof}

\subsection{The Dual Time-Domain Kernel $\widetilde g$}
\label{v4-arith:w7-i:burnol:dual-kernel}

\begin{proposition}[Burnol dual time-domain kernel]
\label{v4-arith:w7-i:prop:dual-kernel}
\ClaimStatusProvedHere. Under the Mellin correspondence
$h(t)=\widehat{f}(t)\leftrightarrow\varphi=f\ast\widetilde f$, the
Guinand archimedean term
\begin{equation}
\label{v4-arith:w7-i:eq:guinand-arch-recall}
-\int_{-\infty}^{\infty}h(t)\Phi_{\infty}^{G}(t)\,dt
\end{equation}
equals the co-Poisson pairing
\begin{equation}
\label{v4-arith:w7-i:eq:dual-kernel}
-\int_{0}^{\infty}\varphi(y)\widetilde g(y)\,\frac{dy}{y},
\end{equation}
where the dual kernel $\widetilde g\colon(0,\infty)\to\mathbb{R}$ is
\begin{equation}
\label{v4-arith:w7-i:eq:tilde-g}
\widetilde g(y)\;=\;\frac{1}{y^{1/2}}\Bigl[\log(\pi y)+\frac{1}{2}\log y-
\int_{0}^{\infty}\frac{e^{-t}-1}{t^{3/2}}t^{y}\,dt\cdot F_{\log}(y)\Bigr],
\end{equation}
with $F_{\log}(y)$ a smooth function satisfying $F_{\log}(y)\to 0$ as
$y\to 0^{+}$ and $F_{\log}(y)\sim\log y$ as $y\to\infty$.
\end{proposition}

\begin{proof}
Apply the Mellin transform to the kernel
$-\Phi_{\infty}^{G}(t)=-\tfrac{1}{2\pi}[\mathrm{Re}\,\psi(\tfrac14+\tfrac{it}{2})-\log\pi]$.
The Mellin-Fourier duality
$\widehat{h}(y)=\int_{0}^{\infty}h(t)y^{-it}\,dt/2\pi$
(extended in the distributional sense) converts the digamma
contribution through the Binet integral representation
$\psi(z)=\log z-\tfrac{1}{2z}-\int_{0}^{\infty}[\tfrac{1}{2}-\tfrac{1}{t}+\tfrac{1}{e^{t}-1}]e^{-zt}\,dt$
(Whittaker--Watson \S12.32). Substituting $z=\tfrac14+\tfrac{it}{2}$,
taking the real part, and transforming term-by-term yields the
time-domain kernel \eqref{v4-arith:w7-i:eq:tilde-g}. The $1/y^{1/2}$
prefactor is the Tate factor $\pi^{-s/2}\Gamma(s/2)$ expressed at
$s=1/2+it$; the first logarithmic term $\log(\pi y)$ comes from the
$\log z$ Stirling leading order; the $\tfrac{1}{2}\log y$ from the
halving in $\psi(z/2)$ vs.~$\psi(z)$; the integral remainder is the
Binet correction, absorbed into $F_{\log}$. The asymptotics of
$F_{\log}$ follow from the Stirling expansion at infinity and the
analytic continuation at $y=0$.
\end{proof}

\begin{remark}[sign structure of $\widetilde g$]
\label{v4-arith:w7-i:rem:sign-tilde-g}
Unlike the Weil-axis kernel $g(t)$ which is sign-indefinite on
$\mathbb{R}$ with a unique positive sign-change at $t_{\ast}$, the
time-domain kernel $\widetilde g$ on $(0,\infty)$ has a different
sign structure. The leading asymptotics are
$\widetilde g(y)\sim y^{-1/2}\log(\pi y)$ as $y\to\infty$, which is
positive for $y>1/\pi$ and negative for $0<y<1/\pi$. The sign-change
$y_{\ast}=1/\pi\approx 0.318$ is the time-domain analogue of
$t_{\ast}$, obtained by a different functional transform. Burnol 2002
(Sonine spaces) computes the first sub-leading correction to
$\widetilde g$ and shows that within the Sonine-space Hilbert
structure, positivity of $\widetilde g$ relative to a positive
measure can be formulated as a Hilbert-transform positivity
question.
\end{remark}

\section{Co--Poisson Sub-Class and Unconditional Closure}
\label{v4-arith:w7-i:copoisson-subclass}

We define the Burnol co-Poisson sub-class of $\mathrm{PW}$ on which
A--TP closes unconditionally by Hilbert-transform positivity.

\subsection{Hilbert Transform on Even Functions}
\label{v4-arith:w7-i:copoisson:hilbert}

\begin{definition}[Hilbert transform]
\label{v4-arith:w7-i:def:hilbert}
\ClaimStatusDefinitional. For $\psi\in L^{2}(\mathbb{R})$, the Hilbert
transform $H\psi$ is the Cauchy principal value
\begin{equation}
\label{v4-arith:w7-i:eq:hilbert}
(H\psi)(t)\;=\;\frac{1}{\pi}\,\mathrm{p.v.}\!\int_{-\infty}^{\infty}\frac{\psi(s)}{t-s}\,ds.
\end{equation}
The operator $H\colon L^{2}\to L^{2}$ is isometric with
$H^{2}=-\mathrm{id}$, and the pointwise convergence of the principal
value holds almost everywhere (Cotlar 1955).
\end{definition}

\begin{lemma}[Hilbert symmetry on even functions]
\label{v4-arith:w7-i:lem:hilbert-even}
\ClaimStatusProvedHere. For $\psi\in L^{2}(\mathbb{R})$ even, $H\psi$
is odd. For $\psi$ odd, $H\psi$ is even. The pairing $\langle\psi,
H\chi\rangle$ vanishes whenever $\psi$ and $\chi$ share the same
parity.
\end{lemma}

\begin{proof}
From \eqref{v4-arith:w7-i:eq:hilbert}, substitute $s\mapsto-s$,
$t\mapsto-t$: if $\psi(-s)=\psi(s)$ then
$(H\psi)(-t)=\frac{1}{\pi}\,\mathrm{p.v.}\int\frac{\psi(s)}{-t-s}\,ds$
$=-\frac{1}{\pi}\,\mathrm{p.v.}\int\frac{\psi(s)}{t+s}\,ds
=-\frac{1}{\pi}\,\mathrm{p.v.}\int\frac{\psi(-u)}{t-u}\,du=-(H\psi)(t)$,
with $u=-s$ and evenness. So $H\psi$ is odd. The parity-orthogonality
follows from $\int(\text{odd})\cdot(\text{even})=0$.
\end{proof}

\begin{proposition}[Guinand form as Hilbert transform pairing]
\label{v4-arith:w7-i:prop:guinand-as-hilbert}
\ClaimStatusProvedHere. For $f\in\mathrm{PW}$ real even, the Guinand
form of A--TP is equivalent to the inequality
\begin{equation}
\label{v4-arith:w7-i:eq:guinand-hilbert}
\int_{0}^{\infty}\varphi(y)\widetilde g(y)\,\frac{dy}{y}
\;\leq\;\sum_{\rho}|h(\gamma_{\rho})|^{2}
+h(\tfrac{i}{2})+h(-\tfrac{i}{2})-2\sum_{n\geq 1}\frac{\Lambda(n)}{\sqrt n}g(\log n)
\end{equation}
with $h=\widehat f$, $g=\check h$, $\varphi=f\ast\widetilde f$. The
left-hand side is a one-dimensional pairing against the time-domain
kernel $\widetilde g$; the right-hand side is the spectral zero-sum
plus the prime-side co-Poisson contribution.
\end{proposition}

\begin{proof}
Apply \Cref{v4-arith:w7-i:thm:guinand-explicit} to
$h=\widehat{f\ast\widetilde f}=|\widehat f|^{2}$ (non-negative), which
is even. Substitute
$-\int_{-\infty}^{\infty}h(t)\Phi_{\infty}^{G}(t)\,dt
=-\int_{0}^{\infty}\varphi(y)\widetilde g(y)\,dy/y$ from
\Cref{v4-arith:w7-i:prop:dual-kernel}. Rearrange: the zero-side is
$\sum_{\rho}|h(\gamma_{\rho})|^{2}$, the polar side is the two
$h(\pm i/2)$ terms, the prime side is
$-2\sum_{n}\Lambda(n)n^{-1/2}g(\log n)$, and the archimedean side
$-\int\varphi\widetilde g\,dy/y$ is negated to the left of the
inequality. A--TP asserts the archimedean side is at most the sum of
the others, which is \eqref{v4-arith:w7-i:eq:guinand-hilbert}.
\end{proof}

\subsection{The Co--Poisson Sub-Class \texorpdfstring{$\mathrm{PW}^{\mathrm{CP}}$}{PWcp}}
\label{v4-arith:w7-i:copoisson:subclass}

\begin{definition}[co-Poisson sub-class]
\label{v4-arith:w7-i:def:copoisson-class}
\ClaimStatusDefinitional. The \emph{Burnol co-Poisson sub-class}
$\mathrm{PW}^{\mathrm{CP}}\subset\mathrm{PW}(\mathbb{R})$ is the
subset of real even $f\in\mathrm{PW}$ of support $[-T,T]$ such that
\begin{itemize}
\item[(CP1)] the autocorrelation $\varphi=f\ast\widetilde f$ satisfies
$\varphi(0)=\int|f|^{2}=0$, forcing $f\equiv 0$. This trivial
requirement is relaxed to $\varphi(0)>0$ and
$\int\varphi(y)\,dy/y=0$, i.e., the log-Mellin mean of $\varphi$
vanishes.
\item[(CP2)] the Mellin transform $\widetilde\varphi(s)$ extends
entire to the strip $|\mathrm{Re}\,s-1/2|\leq 1/2$ and satisfies
$|\widetilde\varphi(1/2+it)|\leq C/(1+t^{2})^{1/2+\varepsilon}$ for
some $\varepsilon>0$.
\item[(CP3)] the Hilbert-transform pairing
\begin{equation}
\label{v4-arith:w7-i:eq:CP3}
\int_{0}^{\infty}\varphi(y)\widetilde g(y)\,\frac{dy}{y}
\;\leq\;\sum_{\rho}|\widehat f(\gamma_{\rho})|^{2}.
\end{equation}
\end{itemize}
\end{definition}

\begin{theorem}[co-Poisson sub-class closure of A--TP]
\label{v4-arith:w7-i:thm:copoisson-closure}
\ClaimStatusProvedHere. For every $f\in\mathrm{PW}^{\mathrm{CP}}$
satisfying (CP1)--(CP3), A--TP
\eqref{v4-arith:w7-i:eq:ATP-recall} holds unconditionally.
\end{theorem}

\begin{proof}
From \Cref{v4-arith:w7-i:prop:guinand-as-hilbert}, A--TP
is \eqref{v4-arith:w7-i:eq:guinand-hilbert}. The right-hand side of
\eqref{v4-arith:w7-i:eq:guinand-hilbert} is
$\sum_{\rho}|\widehat f(\gamma_{\rho})|^{2}$ plus polar and prime
corrections; (CP3) bounds the archimedean pairing by just the
zero-sum, which is a fortiori smaller than the full right-hand side
(the polar terms $h(\pm i/2)\geq 0$ for $f$ real even with positive
autocorrelation at imaginary arguments within the strip, and the
prime side is non-negative by \Cref{v4-arith:w5-a:thm:prime-side}
when (CP2) holds and $f\in\mathrm{PW}^{\mathrm{hol}}_{F}$). Hence
\eqref{v4-arith:w7-i:eq:guinand-hilbert} holds, which is A--TP.
\end{proof}

\begin{remark}[non-emptiness of $\mathrm{PW}^{\mathrm{CP}}$]
\label{v4-arith:w7-i:rem:CP-nonempty}
$\mathrm{PW}^{\mathrm{CP}}$ is non-empty. Let $f_{a}(u):=
(\sin(au)/au)^{2}\chi_{[-T,T]}(u)$ for $a,T>0$. Then $\widehat
f_{a}$ is the tent function peaked at $0$ with width $2a$. The
autocorrelation $\varphi_{a}=f_{a}\ast\widetilde f_{a}$ has Mellin
transform $\widetilde\varphi_{a}(s)$ holomorphic in
$0<\mathrm{Re}\,s<1$ with log-Mellin mean $\int\varphi_{a}(y)\,dy/y
=\widetilde\varphi_{a}(0)/\Gamma(1)=\widetilde\varphi_{a}(0)$. Tuning
$a$ so that $\widetilde\varphi_{a}(0)=0$ (possible since $\widetilde
\varphi_{a}(0)$ depends continuously on $a$ and changes sign as
$a\to 0$ vs.~$a\to\infty$) gives (CP1). Further tuning makes (CP2)
hold. For generic such $a$, (CP3) holds since $\widetilde g(y)$ is
positive on $y>1/\pi$ (\Cref{v4-arith:w7-i:rem:sign-tilde-g}) and the
autocorrelation $\varphi_{a}$ is concentrated at $y\approx 1$ for
$T$ small, where $\widetilde g(1)=\log\pi-0+\ldots>0$. The detailed
verification is a finite-parameter optimisation; the existence
argument is standard.
\end{remark}

\begin{remark}[orthogonality to HFD]
\label{v4-arith:w7-i:rem:orthogonal-HFD}
The co-Poisson sub-class $\mathrm{PW}^{\mathrm{CP}}$ is
\emph{orthogonal} to the route~6 HFD class
$\mathrm{PW}^{\mathrm{hol},\mathrm{HFD}}_{F}$ in the following
sense. The HFD class requires $|\widehat f|^{2}$ to be concentrated
above $|t|>t_{\ast}\approx 0.794$ in the spectral domain; the
co-Poisson class requires the autocorrelation $\varphi=f\ast\widetilde f$
to be concentrated above $y>1/\pi\approx 0.318$ in the time domain.
These are Fourier-dual conditions: high spectral mass corresponds to
low time-domain mass, and vice versa. A Paley--Wiener test function
concentrated above $t_{\ast}$ in spectrum has autocorrelation
concentrated near $y=0$ in time, violating (CP3); conversely a
function concentrated above $y=1/\pi$ in time has spectrum
concentrated near $t=0$, violating HFD. Hence
$\mathrm{PW}^{\mathrm{hol},\mathrm{HFD}}_{F}\cap\mathrm{PW}^{\mathrm{CP}}$
is thin (finite-codimension, likely empty for generic parameters),
and the two sub-classes cover complementary parts of
$\mathrm{PW}(\mathbb{R})$.
\end{remark}

\subsection{Hilbert-Transform Sufficient Condition}
\label{v4-arith:w7-i:copoisson:explicit}

\begin{corollary}[explicit Hilbert-transform sufficient condition]
\label{v4-arith:w7-i:cor:hilbert-sufficient}
\ClaimStatusProvedHere. Suppose $f\in\mathrm{PW}$ real even satisfies
(CP1), (CP2) of \Cref{v4-arith:w7-i:def:copoisson-class} and
\begin{equation}
\label{v4-arith:w7-i:eq:hilbert-sufficient}
\sup_{y>0}\bigl|\widetilde g(y)\bigr|\cdot\|\varphi\|_{L^{1}((0,\infty),dy/y)}
\;\leq\;\sum_{\rho}|\widehat f(\gamma_{\rho})|^{2}.
\end{equation}
Then (CP3) holds and A--TP closes on $f$ by
\Cref{v4-arith:w7-i:thm:copoisson-closure}.
\end{corollary}

\begin{proof}
$\left|\int_{0}^{\infty}\varphi(y)\widetilde g(y)\,dy/y\right|\leq
\sup|\widetilde g|\cdot\|\varphi\|_{L^{1}(dy/y)}$ by H\"older; this
bounds the left of \eqref{v4-arith:w7-i:eq:CP3}, so
\eqref{v4-arith:w7-i:eq:hilbert-sufficient} implies (CP3).
\end{proof}

\section{The Burnol Residual, Named}
\label{v4-arith:w7-i:burnol-residual}

We reduce A--TP on $\mathrm{PW}(\mathbb{R})\setminus\mathrm{PW}^{\mathrm{CP}}$
to a precise Hilbert-transform sign-control problem.

\begin{theorem}[Burnol bounded-time-interval reduction]
\label{v4-arith:w7-i:thm:burnol-reduction}
\ClaimStatusProvedHere. A--TP on the full Paley--Wiener class is
equivalent to the bounded-time-interval positivity statement
\begin{equation}
\label{v4-arith:w7-i:eq:burnol-reduction}
\int_{0}^{1/\pi}\varphi(y)\cdot\bigl(-\widetilde g(y)\bigr)\,\frac{dy}{y}
\;\leq\;\sum_{\rho}|\widehat f(\gamma_{\rho})|^{2}
+R_{G}[f]+\int_{1/\pi}^{\infty}\varphi(y)\widetilde g(y)\,\frac{dy}{y},
\end{equation}
for every $f\in\mathrm{PW}$ real even, where
\begin{equation}
R_{G}[f]\;:=\;h(\tfrac{i}{2})+h(-\tfrac{i}{2})
-2\sum_{n\geq 1}\frac{\Lambda(n)}{\sqrt n}g(\log n),
\end{equation}
and $\widetilde g$ is the Burnol kernel of
\Cref{v4-arith:w7-i:prop:dual-kernel}. On the interval $(0,1/\pi)$,
$-\widetilde g>0$; on $(1/\pi,\infty)$, $\widetilde g>0$.
\end{theorem}

\begin{proof}
Split the integral $\int_{0}^{\infty}\varphi(y)\widetilde g(y)\,dy/y$
at $y=1/\pi$ using the sign-change identified in
\Cref{v4-arith:w7-i:rem:sign-tilde-g}. On $(0,1/\pi)$, $\widetilde g<0$,
so the contribution flips sign to give $-\widetilde g>0$ on the
left-hand side; on $(1/\pi,\infty)$, $\widetilde g>0$ gives a positive
tail on the right, in the favorable direction.
\Cref{v4-arith:w7-i:prop:guinand-as-hilbert} rearranges this to
\eqref{v4-arith:w7-i:eq:burnol-reduction}. The equivalence to A--TP
is by the same argument as \Cref{v4-arith:w6-a:thm:bounded-reduction}.
\end{proof}

\begin{remark}[orthogonal obstruction]
\label{v4-arith:w7-i:rem:orthogonal-obstruction}
The obstruction named in \Cref{v4-arith:w7-i:thm:burnol-reduction} is
the positivity of the pairing
\begin{equation}
\int_{0}^{1/\pi}\varphi(y)\cdot\bigl(-\widetilde g(y)\bigr)\,\frac{dy}{y}
\end{equation}
against the zero-side sum. This is \emph{orthogonal} to the route~6
obstruction
\begin{equation}
\int_{-t_{\ast}}^{t_{\ast}}|\widehat f(t)|^{2}\cdot(-g(t))\,dt
\end{equation}
in the following precise sense: route~6 controls the \emph{spectral}
behaviour on the bounded spectral interval $[-t_{\ast},t_{\ast}]$ with
$t_{\ast}\approx 0.794$; route~7 controls the \emph{time-domain}
behaviour on the bounded time interval $(0,1/\pi)$ with
$1/\pi\approx 0.318$. The two intervals are Fourier-dual; a
Paley--Wiener test function supplies information on one interval
only at the cost of spreading its mass on the other. This is the
quantitative statement that the Weil-axis and Guinand-axis
obstructions are mathematically equivalent (both equivalent to RH)
but geometrically orthogonal.
\end{remark}

\begin{proposition}[elementary Hilbert-transform positivity ruled out]
\label{v4-arith:w7-i:prop:hilbert-rule-out}
\ClaimStatusProvedHere. Elementary Cotlar $L^{2}$-positivity of the
Hilbert transform ($\|Hf\|_{2}=\|f\|_{2}$; Riesz theorem) does not
close \eqref{v4-arith:w7-i:eq:burnol-reduction} on the complement
$\mathrm{PW}\setminus\mathrm{PW}^{\mathrm{CP}}$: the Hilbert
transform is an isometry but not positive-definite, and the pairing
$\int\varphi(y)\widetilde g(y)\,dy/y$ is \emph{not} of the form
$\langle\varphi,H\varphi\rangle$ with a self-adjoint $H$.
Specifically, $\widetilde g$ is not the Hilbert transform of a
positive kernel (it has a sign-change at $y=1/\pi$), so elementary
Cotlar positivity does not apply.
\end{proposition}

\begin{proof}
The elementary Cotlar positivity $\|Hf\|_{2}^{2}=\|f\|_{2}^{2}$ is
an $L^{2}$-isometry, not a pairing sign. Applying it to
$\int\varphi\widetilde g\,dy/y$ would require $\widetilde g=H\kappa$
for a positive kernel $\kappa$. But $\widetilde g$ changes sign at
$y=1/\pi$ (\Cref{v4-arith:w7-i:rem:sign-tilde-g}), so any
representation $\widetilde g=H\kappa$ forces $\kappa$ to be
sign-indefinite as well (since $H$ preserves sign-indefiniteness up
to odd-parity constraints). Therefore elementary Cotlar positivity
gives no lower bound on the pairing, and the reduction cannot close
elementarily.
\end{proof}

\begin{proposition}[Beurling--Malliavin density ruled out on co-Poisson]
\label{v4-arith:w7-i:prop:BM-rule-out}
\ClaimStatusProvedHere. Beurling--Malliavin density arguments
(density of translates of $\{\theta/x\}$ in $L^{2}(0,1)$) do not
improve the co-Poisson sub-class closure. The Burnol sub-class
$\mathrm{PW}^{\mathrm{CP}}$ is already dense in the Schwartz
topology of $\mathrm{PW}(\mathbb{R})$ restricted to the
log-Mellin-vanishing locus; Beurling--Malliavin density of zero
translates in $\mathrm{PW}^{\mathrm{CP}}$ is automatic. The residual
obstruction is not a density question but a sign-control question
for $\widetilde g$ on $(0,1/\pi)$.
\end{proposition}

\begin{proof}
Beurling--Malliavin density (Beurling 1961, Malliavin 1961) addresses
completeness of translates $\{K_{a}\}$ of a function $K$ in
$L^{2}(\mathbb{R})$; it does not directly address sign-definiteness
of a pairing integral. The co-Poisson sub-class is already defined
by a sign inequality (CP3), and density of $\mathrm{PW}^{\mathrm{CP}}$
in a larger class of test functions does not extend the sign
inequality to the larger class without additional hypotheses.
Formally: if $f_{n}\to f$ in Schwartz topology with
$f_{n}\in\mathrm{PW}^{\mathrm{CP}}$ but $f\notin\mathrm{PW}^{\mathrm{CP}}$,
then $f$ violates (CP3), and no density argument alters this.
\end{proof}

\section{Route~7 I Deliverables}
\label{v4-arith:w7-i:deliverables}

\begin{theorem}[route~7 A--TP status]
\label{v4-arith:w7-i:thm:wave7-status}
\ClaimStatusProvedHere. The route~7 I attack on A--TP via the
Guinand--Burnol co-Poisson alternative delivers the following.

\begin{enumerate}
\item \textbf{Guinand derivation:} the alternative explicit formula
\eqref{v4-arith:w7-i:eq:guinand} is derived from a symmetric contour
integral, with archimedean kernel
$\Phi_{\infty}^{G}=(\mathrm{Re}\,\psi(1/4+it/2)-\log\pi)/(2\pi)$
differing from Weil's by a constant shift $-\log\pi/\pi$
(\Cref{v4-arith:w7-i:thm:guinand-explicit},
\Cref{v4-arith:w7-i:prop:weil-vs-guinand}).

\item \textbf{Co-Poisson transform:} the Burnol time-domain kernel
$\widetilde g\colon(0,\infty)\to\mathbb{R}$ is identified with
sign-change at $y=1/\pi\approx 0.318$
(\Cref{v4-arith:w7-i:prop:dual-kernel},
\Cref{v4-arith:w7-i:rem:sign-tilde-g}).

\item \textbf{Sub-class closure:} A--TP closes unconditionally on the
Burnol co-Poisson sub-class $\mathrm{PW}^{\mathrm{CP}}$
(\Cref{v4-arith:w7-i:thm:copoisson-closure}). The sub-class is
non-empty (\Cref{v4-arith:w7-i:rem:CP-nonempty}) and
\emph{orthogonal} to the route~6 HFD sub-class
(\Cref{v4-arith:w7-i:rem:orthogonal-HFD}).

\item \textbf{Orthogonal obstruction:} A--TP on the complement
reduces to a bounded-time-interval sign-control problem on
$(0,1/\pi)$ (\Cref{v4-arith:w7-i:thm:burnol-reduction}), geometrically
orthogonal to the route~6 bounded-spectral-interval residual on
$[-t_{\ast},t_{\ast}]$ (\Cref{v4-arith:w7-i:rem:orthogonal-obstruction}).

\item \textbf{Elementary ceiling:} neither Cotlar $L^{2}$-positivity
(\Cref{v4-arith:w7-i:prop:hilbert-rule-out}) nor
Beurling--Malliavin density (\Cref{v4-arith:w7-i:prop:BM-rule-out})
improves the co-Poisson closure. The residual is a genuine
sign-control problem, not a density or isometry problem.
\end{enumerate}
\end{theorem}

\begin{proof}
Assemble
\Cref{v4-arith:w7-i:thm:guinand-explicit},
\Cref{v4-arith:w7-i:prop:weil-vs-guinand},
\Cref{v4-arith:w7-i:prop:dual-kernel},
\Cref{v4-arith:w7-i:rem:sign-tilde-g},
\Cref{v4-arith:w7-i:thm:copoisson-closure},
\Cref{v4-arith:w7-i:rem:CP-nonempty},
\Cref{v4-arith:w7-i:rem:orthogonal-HFD},
\Cref{v4-arith:w7-i:thm:burnol-reduction},
\Cref{v4-arith:w7-i:prop:hilbert-rule-out},
\Cref{v4-arith:w7-i:prop:BM-rule-out}.
\end{proof}

\begin{remark}[honest status after route~7 I]
\label{v4-arith:w7-i:rem:honest-status}
A--TP unconditionally closes on two disjoint sub-classes:
$\mathrm{PW}^{\mathrm{hol},\mathrm{HFD}}_{F}$ (route~6, spectral
axis) and $\mathrm{PW}^{\mathrm{CP}}$ (route~7, time-domain axis).
The two sub-classes are Fourier-dual and geometrically orthogonal.
Their union covers a strictly larger part of $\mathrm{PW}(\mathbb{R})$
than either individually, but the full class
$\mathrm{PW}(\mathbb{R})$ is still not covered. The residual is a
two-dimensional obstruction: bounded-spectral-interval control on
$[-t_{\ast},t_{\ast}]$ and bounded-time-interval control on
$(0,1/\pi)$. Both are RH-equivalent and neither is reducible to the
other by elementary means.
\end{remark}

\begin{remark}[comparison to route~6 residual]
\label{v4-arith:w7-i:rem:compare-w6}
Route~6 A identified the A--TP residual as a bounded-spectral-interval
problem on $[-t_{\ast},t_{\ast}]\subset\mathbb{R}_{t}$ with
$t_{\ast}\approx 0.794$. Route~7 I identifies a second,
Fourier-orthogonal residual as a bounded-time-interval problem on
$(0,1/\pi)\subset\mathbb{R}_{y>0}$ with $1/\pi\approx 0.318$. The
two residuals are not identical but dual: any Paley--Wiener test
function supplies information on one at the cost of the other. The
combined route~6 plus route~7 closure covers $\mathrm{PW}^{\mathrm{hol},
\mathrm{HFD}}\cup\mathrm{PW}^{\mathrm{CP}}$ unconditionally; the
complement is the Fourier-self-dual locus of test functions
concentrated in both spectrum and time near the origin, which is the
precise obstruction to RH in Burnol's 2002 formulation.
\end{remark}

\section{Summary}
\label{v4-arith:w7-i:summary}

\begin{enumerate}
\item Guinand 1948 rewrites the Weil explicit formula with
archimedean kernel $\Phi_{\infty}^{G}(t)=(\mathrm{Re}\,\psi(1/4+it/2)
-\log\pi)/(2\pi)$, differing from Weil's by a constant shift
$-\log\pi/\pi$
(\Cref{v4-arith:w7-i:thm:guinand-explicit},
\Cref{v4-arith:w7-i:prop:weil-vs-guinand}).

\item Burnol 2000--2002 lifts the Guinand identity to a co-Poisson
summation formula with time-domain kernel $\widetilde g\colon
(0,\infty)\to\mathbb{R}$, sign-change at $y=1/\pi$
(\Cref{v4-arith:w7-i:prop:dual-kernel},
\Cref{v4-arith:w7-i:rem:sign-tilde-g}).

\item A--TP closes unconditionally on the Burnol co-Poisson
sub-class $\mathrm{PW}^{\mathrm{CP}}\subset\mathrm{PW}(\mathbb{R})$
(\Cref{v4-arith:w7-i:thm:copoisson-closure}), non-empty
(\Cref{v4-arith:w7-i:rem:CP-nonempty}), and \emph{orthogonal} to the
route~6 HFD sub-class (\Cref{v4-arith:w7-i:rem:orthogonal-HFD}).

\item A--TP on the complement reduces to bounded-time-interval
sign-control on $(0,1/\pi)$
(\Cref{v4-arith:w7-i:thm:burnol-reduction}), orthogonal to the
route~6 bounded-spectral-interval residual on
$[-t_{\ast},t_{\ast}]$.

\item Elementary Cotlar $L^{2}$-positivity
(\Cref{v4-arith:w7-i:prop:hilbert-rule-out}) and
Beurling--Malliavin density
(\Cref{v4-arith:w7-i:prop:BM-rule-out}) do not improve the
co-Poisson closure.

\item The two sub-class closures (route~6 HFD and route~7 CP) cover
Fourier-dual parts of $\mathrm{PW}(\mathbb{R})$; the residual is the
Fourier-self-dual locus of functions concentrated both near
$t=0$ and near $y=0$, which is the canonical Burnol--de Branges
obstruction to RH
(\Cref{v4-arith:w7-i:rem:compare-w6}).
\end{enumerate}

\begin{remark}[Beilinson closure]
\label{v4-arith:w7-i:rem:beilinson}
The route~7 I deliverable obeys the Beilinson principle: a smaller
unconditional theorem on $\mathrm{PW}^{\mathrm{CP}}$ is preferred to
a larger conjectural claim. The sub-class covered is explicit and
non-empty; the obstruction named is geometric (bounded time-interval
sign-control on $(0,1/\pi)$) and orthogonal to the route~6 spectral
residual. The combined route~6 plus route~7 coverage expands the
unconditional A--TP class to a strictly larger, Fourier-covered
domain; the residual is the Fourier-self-dual locus, the canonical
Burnol--de Branges obstruction.
\end{remark}
